\documentclass[11pt]{article}

\usepackage[margin=1in]{geometry}
\usepackage{amsmath,amssymb,amsthm}
\usepackage{enumitem}
\usepackage{hyperref}
\usepackage{xcolor}
\usepackage{booktabs}

\hypersetup{colorlinks=true, linkcolor=blue, urlcolor=blue, citecolor=blue}

\title{TOAE209/TOAH209 -- Design and Analysis of Algorithms\\[4pt]
\large Week 7: Theoretical Questions}
\author{Foreign Trade University}
\date{}

\begin{document}
\maketitle

\section*{Instructions}

Part A contains 16 questions requiring written explanations or short proofs. Part B is a
10-question quiz (true/false and multiple choice). Attempt every question on your own before
consulting Part C (Solutions) at the end of this document.

\section*{Part A: Theory Questions}

\begin{enumerate}[label=\textbf{A\arabic*.}, leftmargin=2.2em]

\item State the closest-pair-of-points problem precisely. What is the running time of the
brute-force algorithm, and why is it $\Theta(n^2)$?

\item In the divide-and-conquer closest-pair algorithm, after recursing on the left and right
halves we obtain $\delta = \min(\delta_L, \delta_R)$. Explain why any crossing pair (one point
in each half) that is closer than $\delta$ must lie entirely within the vertical strip of
width $2\delta$ centred on the dividing line.

\item Prove the ``sparsity'' lemma: when the strip points are sorted by $y$-coordinate, two
points within distance $\delta$ of each other are within a constant number of positions in
that order. Use the $\delta/2$-box argument and explain why no box contains two points.

\item Why does the closest-pair algorithm pre-sort the points by $y$-coordinate \emph{once} at
the top level, and pass the $y$-sorted sublists down through the recursion, rather than
re-sorting the strip by $y$ inside each call? What would re-sorting do to the recurrence?

\item Write and solve the recurrence for the divide-and-conquer closest-pair algorithm.
Identify the cost of the combine (strip) step and explain why it is linear.

\item Derive the naive $\Theta(n^2)$ recurrence for integer multiplication by splitting each
$n$-digit number into two halves. Exactly which four products appear, and why does the master
theorem give $\Theta(n^2)$?

\item State Karatsuba's identity: how are the three products $z_0, z_1, z_2$ defined, and why
does $z_1$ equal the desired middle coefficient $x_1 y_0 + x_0 y_1$? Write the resulting
recurrence and solve it.

\item Strassen's algorithm multiplies two $n \times n$ matrices using $7$ recursive
multiplications of $(n/2) \times (n/2)$ blocks instead of $8$. Write both recurrences ($8$-way
and $7$-way) with an $O(n^2)$ combine step and solve each. For which $n$ does Strassen perform
strictly fewer scalar multiplications?

\item Explain why multiplying two polynomials given by their coefficient vectors is exactly
the convolution of those vectors. Give the formula for the $k$-th coefficient of the product.

\item A degree-$(n-1)$ polynomial can be represented by $n$ coefficients or by its values at
$n$ distinct points. Explain why multiplication is $O(n)$ in the value representation but
$\Theta(n^2)$ in the coefficient representation, and why this motivates the FFT.

\item Define the $n$-th roots of unity $\omega_n^k$. State and prove the squaring (halving)
property $(\omega_n^k)^2 = \omega_{n/2}^k$ and the cancellation property
$\omega_n^{k + n/2} = -\omega_n^k$.

\item Derive the even/odd splitting identity $A(x) = A_{\text{even}}(x^2) + x\,
A_{\text{odd}}(x^2)$ and explain how, combined with the two root-of-unity properties, it lets
the FFT compute the length-$n$ DFT from two length-$(n/2)$ DFTs in $O(n)$ extra work. Write
the resulting recurrence.

\item Show that the inverse DFT has the same form as the forward DFT with $\omega_n$ replaced
by $\omega_n^{-1}$ and an extra factor of $1/n$. Prove the key orthogonality identity
$\frac{1}{n}\sum_{k=0}^{n-1}\omega_n^{k(j-j')} = [\,j = j'\,]$.

\item Describe the complete $O(n \log n)$ algorithm for multiplying two polynomials via the
FFT (padding, forward transforms, pointwise product, inverse transform, rounding). Prove its
correctness, i.e.\ that $\hat C[k] = C(\omega^k)$ for the true product $C = AB$.

\item Explain how counting the number of contiguous subarrays whose sum lies in $[\ell, u]$
reduces to a problem on the prefix-sum array, and how a mergesort with a two-pointer count
during the merge solves it in $O(n \log n)$.

\item Describe the $O(\log(m+n))$ algorithm for the median of two sorted arrays. What is the
partition invariant that the binary search maintains, and why does it guarantee the median is
found?

\end{enumerate}

\newpage
\section*{Part B: Quiz}

\begin{enumerate}[label=\textbf{B\arabic*.}, leftmargin=2.2em]

\item \textbf{(Multiple choice)} The divide-and-conquer closest-pair algorithm runs in:
\begin{enumerate}[label=(\alph*)]
    \item $\Theta(n^2)$
    \item $\Theta(n \log n)$
    \item $\Theta(n)$
    \item $\Theta(n \log^2 n)$
\end{enumerate}

\item \textbf{(True/False)} In the closest-pair strip, each point need only be compared with a
constant number of its neighbours in $y$-sorted order (e.g.\ the next 7), regardless of $n$.

\item \textbf{(Multiple choice)} Karatsuba's algorithm multiplies two $n$-digit integers in
time:
\begin{enumerate}[label=(\alph*)]
    \item $\Theta(n)$
    \item $\Theta(n \log n)$
    \item $\Theta(n^{\log_2 3}) \approx \Theta(n^{1.585})$
    \item $\Theta(n^2)$
\end{enumerate}

\item \textbf{(Short answer)} How many recursive multiplications does Karatsuba use per level,
and how many does the naive split use?

\item \textbf{(Multiple choice)} Strassen's matrix multiplication has running time:
\begin{enumerate}[label=(\alph*)]
    \item $\Theta(n^2)$
    \item $\Theta(n^{\log_2 7}) \approx \Theta(n^{2.807})$
    \item $\Theta(n^3)$
    \item $\Theta(n^2 \log n)$
\end{enumerate}

\item \textbf{(True/False)} Multiplying two polynomials given as coefficient vectors is the
same as convolving those vectors.

\item \textbf{(Short answer)} What is $\omega_n^{k + n/2}$ in terms of $\omega_n^k$, and what
is this property called?

\item \textbf{(Multiple choice)} The recursive radix-2 FFT computes the DFT of a length-$n$
vector (n a power of two) in time:
\begin{enumerate}[label=(\alph*)]
    \item $\Theta(n^2)$
    \item $\Theta(n)$
    \item $\Theta(n \log n)$
    \item $\Theta(\log n)$
\end{enumerate}

\item \textbf{(True/False)} The inverse DFT can be computed by the same FFT routine using the
conjugate root $\omega_n^{-1}$ and dividing the result by $n$.

\item \textbf{(Short answer)} In fast modular exponentiation (square-and-multiply), roughly
how many modular multiplications are needed to compute $a^e \bmod m$?

\end{enumerate}

\newpage
\section*{Part C: Solutions}

\subsection*{Part A}

\begin{enumerate}[label=\textbf{A\arabic*.}, leftmargin=2.2em]

\item \textbf{Problem.} Given $n$ points $p_1, \ldots, p_n$ in the plane, find the pair
minimizing the Euclidean distance $d(p_i,p_j) = \sqrt{(x_i-x_j)^2 + (y_i-y_j)^2}$. The
brute-force algorithm computes the distance for all $\binom{n}{2} = \frac{n(n-1)}{2}$ pairs
and keeps the minimum; this is $\Theta(n^2)$ because it does a constant amount of work for
each of $\Theta(n^2)$ pairs.

\item Suppose a crossing pair $(p, q)$ has $p \in L$, $q \in R$, and $d(p, q) < \delta$. Since
the Euclidean distance is at least the difference in any single coordinate, $|x_p - x_q| \le
d(p,q) < \delta$. The dividing line $x = x^*$ lies between $x_p$ (on the left) and $x_q$ (on
the right), so $|x_p - x^*| \le |x_p - x_q| < \delta$ and similarly $|x_q - x^*| < \delta$.
Hence both points are within horizontal distance $\delta$ of the line, i.e.\ inside the strip.

\item Partition the strip (width $2\delta$) into a grid of square boxes of side $\delta/2$.
Each box has diameter $\frac{\delta}{2}\sqrt{2} = \frac{\delta}{\sqrt 2} < \delta$. No box can
contain two points: two points in the same box would be at distance $< \delta$, but every
pair of points lying on the same side of the dividing line is at distance $\ge \delta$
(guaranteed by the recursive calls), and a box of width $\delta/2$ small enough to straddle
the line still has diameter $< \delta$, so two points in it would contradict $\delta$ being
the current minimum. Now if two strip points are within distance $\delta$, their $y$-values
differ by $\le \delta$, so the second lies within a fixed-size block of boxes around the
first (the strip is $4$ columns wide, and a height of $\delta$ above/below spans a constant
number of rows). Since each box holds at most one point, only a constant number ($\le 15$,
and after a more careful count $\le 7$ in $y$-order) of points can be that close. Hence
comparing each point with its next $7$ (or $15$) neighbours in $y$-order suffices.

\item Pre-sorting by $y$ once costs $O(n \log n)$ at the top level only. Each recursive call
then produces its left/right $y$-sorted sublists by a single linear scan of its already
$y$-sorted input (a partition, not a sort), keeping the per-call combine cost $O(n)$. If
instead we re-sorted the strip by $y$ inside every call, each call would cost $O(n \log n)$,
giving the recurrence $T(n) = 2T(n/2) + O(n \log n) = O(n \log^2 n)$ -- a $\log n$ factor
worse. Pre-sorting preserves the $O(n)$ combine and hence the $O(n \log n)$ total.

\item The recurrence is $T(n) = 2\,T(n/2) + O(n)$. The two recursive calls handle the halves;
the $O(n)$ combine splits $P_y$ in linear time and scans the strip doing a constant number of
distance comparisons per strip point (by A3), which is $O(|S|) = O(n)$. By the master theorem
($a = 2$, $b = 2$, $d = 1$, $a = b^d$), $T(n) = \Theta(n \log n)$. (The initial sort is a
one-time $O(n\log n)$ cost that does not change the bound.)

\item Write $x = x_1 B^{n/2} + x_0$ and $y = y_1 B^{n/2} + y_0$. Then $xy = x_1 y_1 B^n +
(x_1 y_0 + x_0 y_1) B^{n/2} + x_0 y_0$, requiring the four products $x_1 y_1$, $x_1 y_0$,
$x_0 y_1$, $x_0 y_0$, each a multiplication of $(n/2)$-digit numbers, plus $O(n)$ additions and
shifts. So $T(n) = 4\,T(n/2) + O(n)$; with $a = 4$, $b = 2$, $d = 1$, we have $a = 4 > b^d =
2$, so $T(n) = \Theta(n^{\log_2 4}) = \Theta(n^2)$.

\item Define $z_2 = x_1 y_1$, $z_0 = x_0 y_0$, and $z_1 = (x_0 + x_1)(y_0 + y_1) - z_2 - z_0$.
Expanding, $(x_0 + x_1)(y_0 + y_1) = x_0 y_0 + x_0 y_1 + x_1 y_0 + x_1 y_1 = z_0 + (x_0 y_1 +
x_1 y_0) + z_2$, so subtracting $z_0$ and $z_2$ leaves $z_1 = x_0 y_1 + x_1 y_0$, exactly the
middle coefficient. Then $xy = z_2 B^n + z_1 B^{n/2} + z_0$ uses only three recursive
multiplications ($z_0$, $z_2$, and the product inside $z_1$), giving $T(n) = 3\,T(n/2) + O(n)
= \Theta(n^{\log_2 3}) \approx \Theta(n^{1.585})$.

\item The standard recursive 2x2-block multiply does $8$ block multiplications: $T(n) = 8\,
T(n/2) + O(n^2)$, with $a = 8$, $b = 2$, $d = 2$, $a = 8 > b^d = 4$, so $T(n) =
\Theta(n^{\log_2 8}) = \Theta(n^3)$. Strassen uses $7$: $T(n) = 7\, T(n/2) + O(n^2)$, with
$a = 7 > b^d = 4$, so $T(n) = \Theta(n^{\log_2 7}) \approx \Theta(n^{2.807})$. Counting only
scalar multiplications with a $1\times 1$ base case costing $1$: the naive scheme uses
$8^{\log_2 n} = n^3$ and Strassen uses $7^{\log_2 n} = n^{\log_2 7}$; Strassen is strictly
fewer for every $n \ge 2$ (a power of two), since $7 < 8$ at each of the $\log_2 n \ge 1$
levels.

\item If $A(x) = \sum_i a_i x^i$ and $B(x) = \sum_j b_j x^j$, then $A(x)B(x) = \sum_k
\bigl(\sum_{i+j=k} a_i b_j\bigr) x^k$. So the $k$-th coefficient of the product is $c_k =
\sum_{i+j=k} a_i b_j$, which is precisely the (linear) convolution of the coefficient vectors
$(a_i)$ and $(b_j)$.

\item In the value representation, if $A$ and $B$ are known at the same set of $\ge \deg A +
\deg B + 1$ points, then $(AB)$ at each point is just the product of the two values --
$O(n)$ pointwise multiplications total. In the coefficient representation, computing the
convolution directly is the double sum over all $i, j$, costing $\Theta(n^2)$. This asymmetry
motivates the FFT: convert coefficients $\to$ values (evaluate), multiply pointwise in $O(n)$,
then convert values $\to$ coefficients (interpolate). The FFT makes both conversions
$O(n \log n)$, so the whole multiplication is $O(n \log n)$.

\item $\omega_n^k = e^{2\pi i k/n}$ are the $n$ equally spaced points on the unit circle with
$z^n = 1$. \emph{Squaring:} $(\omega_n^k)^2 = e^{2\pi i \cdot 2k/n} = e^{2\pi i k/(n/2)} =
\omega_{n/2}^k$. \emph{Cancellation:} $\omega_n^{k + n/2} = e^{2\pi i (k+n/2)/n} = e^{2\pi i
k/n} \cdot e^{\pi i} = \omega_n^k \cdot (-1) = -\omega_n^k$.

\item Separating even- and odd-indexed coefficients, $A(x) = \sum_j a_{2j} x^{2j} + \sum_j
a_{2j+1} x^{2j+1} = A_{\text{even}}(x^2) + x\, A_{\text{odd}}(x^2)$, where $A_{\text{even}},
A_{\text{odd}}$ have degree $< n/2$. Evaluating at $x = \omega_n^k$: by the squaring property
$(\omega_n^k)^2 = \omega_{n/2}^k$, so $A_{\text{even}}$ and $A_{\text{odd}}$ are evaluated at
the $(n/2)$-th roots of unity -- two DFTs of size $n/2$. By cancellation, for $k < n/2$,
$A(\omega_n^k) = E[k] + \omega_n^k O[k]$ and $A(\omega_n^{k+n/2}) = E[k] - \omega_n^k O[k]$,
so all $n$ outputs are produced from the two half-size DFTs with $n/2$ butterflies, i.e.\
$O(n)$ extra work. The recurrence is $T(n) = 2\,T(n/2) + O(n) = O(n\log n)$.

\item The forward DFT is multiplication by $V$ with $V_{kj} = \omega_n^{kj}$. Claim:
$(V^{-1})_{kj} = \frac1n \omega_n^{-kj}$. Compute $(V \cdot \frac1n \overline V)_{j j'} =
\frac1n \sum_{k} \omega_n^{jk} \omega_n^{-k j'} = \frac1n \sum_k \omega_n^{k(j-j')}$. If $j =
j'$, every term is $1$, giving $\frac1n \cdot n = 1$. If $j \ne j'$, set $r = \omega_n^{j-j'}
\ne 1$; then $\sum_{k=0}^{n-1} r^k = \frac{r^n - 1}{r - 1} = 0$ because $r^n = 1$. Hence the
product is the identity, proving $V^{-1} = \frac1n \overline V$: the inverse DFT is the same
sum with $\omega_n$ replaced by $\omega_n^{-1}$, scaled by $1/n$.

\item \textbf{Algorithm.} Pad $A, B$ with zeros to a common length $m$ (a power of two) with
$m \ge \deg A + \deg B + 1$. Compute $\hat A = \text{FFT}(A)$ and $\hat B = \text{FFT}(B)$
(values at the $m$-th roots of unity). Form $\hat C[k] = \hat A[k]\hat B[k]$ for all $k$.
Compute $C = \text{IFFT}(\hat C)$ and round to integers if the inputs were integer-valued.
\textbf{Correctness:} $\hat C[k] = \hat A[k]\hat B[k] = A(\omega^k) B(\omega^k) = C(\omega^k)$,
so $\hat C$ holds the values of the true product $C = AB$ at the $m$ roots of unity. Since
$\deg C = \deg A + \deg B < m$, these $m$ values determine $C$ uniquely, and the inverse FFT
recovers its coefficients exactly. The cost is three FFTs ($O(m \log m)$) plus the $O(m)$
pointwise product, i.e.\ $O(n \log n)$.

\item Let $S[0], \ldots, S[n]$ be the prefix sums, $S[0] = 0$, $S[t] = \sum_{i<t}
\text{nums}[i]$. A subarray spanning indices $(i, j)$ has sum $S[j] - S[i]$, so we must count
pairs $i < j$ with $\ell \le S[j] - S[i] \le u$. Mergesort the array $S$. When merging two
sorted halves (left indices $i$, right indices $j$, with all left positions before all right
positions in the original order), for each left value $S[i]$ slide two pointers over the
sorted right half to count the $j$ with $S[j] - S[i] \in [\ell, u]$ -- a $\ell \le S[j]-S[i]$
lower pointer and a $S[j]-S[i] \le u$ upper pointer, each advancing monotonically. This counts
all straddling pairs in $O(n)$ per merge level, and recursion handles the within-half pairs,
for $O(n \log n)$ total.

\item Assume WLOG $a$ is the shorter array. Binary-search the number $i$ of elements of $a$
placed on the left side of a partition; set $j = \lceil (m+n)/2\rceil - i$ so the left side
holds exactly $\lceil(m+n)/2\rceil$ elements. Using sentinels $a_{\text{left}} = a[i-1]$ (or
$-\infty$), $a_{\text{right}} = a[i]$ (or $+\infty$), and likewise for $b$, the partition is
\emph{correct} when $a_{\text{left}} \le b_{\text{right}}$ and $b_{\text{left}} \le
a_{\text{right}}$: this invariant says every element on the left is $\le$ every element on the
right, so the boundary is the true median split. Then the median is $\max(a_{\text{left}},
b_{\text{left}})$ for odd total length, or the average of that and $\min(a_{\text{right}},
b_{\text{right}})$ for even length. If $a_{\text{left}} > b_{\text{right}}$ we move $i$ left;
otherwise we move $i$ right. The search runs in $O(\log \min(m,n)) = O(\log(m+n))$.

\end{enumerate}

\subsection*{Part B}

\begin{enumerate}[label=\textbf{B\arabic*.}, leftmargin=2.2em]

\item \textbf{(b) $\Theta(n \log n)$.} The recurrence is $T(n) = 2T(n/2) + O(n)$.

\item \textbf{True.} The $\delta/2$-box / sparsity lemma shows only a constant number of
strip points can be within $\delta$ of any given point in $y$-order.

\item \textbf{(c) $\Theta(n^{\log_2 3}) \approx \Theta(n^{1.585})$.} From $T(n) = 3T(n/2) +
O(n)$.

\item Karatsuba uses \textbf{3} recursive multiplications per level; the naive split uses
\textbf{4}.

\item \textbf{(b) $\Theta(n^{\log_2 7}) \approx \Theta(n^{2.807})$.} From $T(n) = 7T(n/2) +
O(n^2)$.

\item \textbf{True.} The product coefficients are exactly the convolution $c_k = \sum_{i+j=k}
a_i b_j$.

\item $\omega_n^{k+n/2} = -\,\omega_n^k$; this is the \textbf{cancellation property} (a root
and its antipode differ only in sign).

\item \textbf{(c) $\Theta(n \log n)$.} From $T(n) = 2T(n/2) + O(n)$ with $n/2$ butterflies per
level.

\item \textbf{True.} $V^{-1} = \frac1n \overline V$, so running the FFT with $\omega_n^{-1}$
and dividing by $n$ inverts the transform.

\item About $\Theta(\log e)$ modular multiplications: one squaring per bit of $e$ plus one
extra multiply per set bit, so at most $2\lfloor\log_2 e\rfloor + 1$.

\end{enumerate}

\end{document}
